\documentclass[11pt]{article}
\usepackage{amsmath}
\title{What pairwise braking certificates retain and what they miss}
\begin{document}
\maketitle

A relative-speed braking calculation can understate required
clearance when a robot cannot accelerate because it is already
at its speed cap. Even exact pairwise calculations need not
compose when two gaps consume the same middle actuator.
Both obstructions have exact, executable boundaries in the
single-file model.

\paragraph{Two robots with independent caps.}
Positions are ordered and the physical gap is to remain in
\([D_s,D_c]\) forever. Motion is forward only, with velocities
\(0\le v_1\le W\), \(0\le v_2\le V\), \(V\le W\), and
acceleration bounds \(a,A\), respectively, all positive.
No new fault, binding wall, or terminal service is included.
Let \(K=a+A\), \(w=(v_1-v_2)_+\), and
\(e=(w-K(V-v_2)/A)_+\).
The minimum distance to the lower clearance limit needed
at this state is
\[
              C=\frac{w^2}{2K}+\frac{A e^2}{2aK}.
\]
The minimum distance to the upper link limit needed is
\(U=(v_2-v_1)_+^2/(2K)\). The controller can preserve both
limits iff the actual margins meet these bounds.

For shrinkage, the left robot can brake as hard as possible and
the right can accelerate as hard as possible, both clamped to their
speed ranges. These velocities minimize the closing speed pointwise;
no other controls can need less clearance to get to the time they
first match. If the right robot does not hit its cap first, integrating
the closing triangle gives \(w^2/(2K)\). If the cap binds, the
closing speed there is \(e>0\); the remaining closing triangle has
slope \(-a\), not \(-K\). Replacing \(e^2/(2K)\) by \(e^2/(2a)\)
gives the formula. While they equalize, the gap shrinks monotonically,
so the upper constraint cannot become tighter. For increasing
separation, the fastest simultaneous braking and acceleration
meet at a common speed between the initial speeds and at most \(V\),
so neither cap blocks the needed intervention. After either
equalization, maintaining the common speed holds the gap invariant.

With \(a=1,A=4,V=1,W=3\), velocities \((1,0)\) and \((2,1)\)
have the same closing speed one. The necessary clearance is \(1/10\)
and \(1/2\) respectively. With margin \(3/10\), a relative-speed
formula alone accepts both while the second is already unrecoverable.

\paragraph{Three robots with a common acceleration cap.}
Suppose the forward velocities initially are \(v_0+2w,v_0+w,v_0\)
for \(v_0\ge0,w>0\). All have acceleration magnitude cap \(a>0\)
and speeds from zero through at least \(v_0+2w\).
Let \(c_1,c_2\) be the current margins above the lower gap limits,
and assume the upper link margins start valid. Put \(m=w^2/(4a)\).
There exists a controller preserving all constraints forever iff
\[
             c_1\ge m,\quad c_2\ge m,\quad c_1+c_2\ge4m.
\]
The first two inequalities are the exact isolated pair constraints.
The last follows by integrating the unavoidable closing of the
outermost pair: their total initial closing speed is \(2w\)
and can be reduced at no more than \(2a\).

To attain, decelerate the leftmost robot at \(a\) and accelerate the
rightmost at \(a\) for \(T=w/a\). Let the middle robot's speed be
\[
        v_0+w+\theta a\min(t,T-t),\qquad -1\le\theta\le1 .
\]
It stays between the outside speeds. Thus both gaps shrink
monotonically and the total loss is \(4m\). The individual losses
are \(2m-\theta m\) and \(2m+\theta m\). A value of \(\theta\) fits
both margins exactly under the three inequalities. Then hold equal
speed. In the replay the pair budgets \(2/5,2/5\) each permit the
isolated maneuver but not the assembly, while \(1/4,3/4\) permits
the assembly. This is a simultaneous input conflict, not incomplete
knowledge of the individual gaps.

\paragraph{Matched barrier comparison.}
For two robots the published common-cap distance barriers of
Wang, Ames and Egerstedt reduce in one dimension to comparing
relative speed with a square root of distance margin.
Setting the relative acceleration bound to \(K\), their
uncapped-relative stopping distance is \(w^2/(2K)\).
The counterexample is to retaining that expression after changing
the independent speed law, not to their theorem under its own
assumptions. Their compositional method also requires simultaneous
input feasibility; the three-robot case shows why its separate
pairwise premises are not alone sufficient. The new formulas
provide an exact comparator for these restricted states.

For the hidden-gap deadline example, the outward braking margin
of the forced common prefix is positive at receipt 0.33 and
already negative at 0.35, even though both positions still
satisfy the static band. At the exact critical receipt the closed
braking bound is tight, and the previously proved trajectory
grazes the link limit. The published strict positive barrier domain
does not certify that equality case; a strictly positive physical
margin would be a stronger law. The replay checks 720 exact
two-robot integrations and the three-robot construction.
This is a bounded kinematic comparison, not a general
barrier-controller benchmark or completion of R6.
\end{document}
